% problem-000070 4 矿物冶金与氧化还原
\section*{4. 矿物冶金与氧化还原}
最简便的处理办法是⽆氧化剂存在的酸溶反应：向含⻩铜矿的矿物中加⼊硫酸，控制浓度和温度，所得体系显蓝⾊且有臭鸡蛋味的⽓体(A)放出（反应1）；为避免⽓体A的产⽣，可采⽤三价铁盐如 $\mathrm { F e _ { 2 } ( S O _ { 4 } ) _ { 3 } }$ 溶液处理⻩铜矿，所得溶液和反应1的产物类似，但有浅⻩⾊固体物质B⽣成（反应2）；浅⻩⾊固体会阻碍⻩铜矿的溶解，因此，办法之⼀是在类似反应1的体系中引⼊硫酸杆菌类微⽣物，同时通⼊空⽓，反应产物⽆⽓体，也⽆⻩⾊浑浊物（反应3），溶液中阴离⼦为C；但是，微⽣物对反应条件要求较为严格（如温度不能过⾼，酸度应适宜等）。采⽤合适的氧化剂如氯酸钾溶液（⾜量），使之在硫酸溶液中与⻩铜矿反应（反应4），是⼀种更有效的处理⽅法。只是，后两种⽅法可能出现副反应（反应5），⽣成⻩钾铁矾(D)。D是⼀种碱式盐（⽆结晶⽔），含两种阳离⼦且⼆者摩尔之⽐为1:3，它会沉积在⻩铜矿上影响其溶解，应尽量避免。5.008 gD和⾜量硫酸钾在硫酸溶液中反应（反应6），得到15.10 g铁钾矾E，E与明矾（相对分⼦质量为474.4 gmol−1）同构。

\subsection*{4-1}
写出A\~C、E的化学式。

\subsection*{4-2}
通过计算，写出D的化学式。

\subsection*{4-3}
写出反应1\~6的离⼦⽅程式（要求系数为最简整数⽐）。
